% !TeX program = xelatex
% 题源均在 27302模拟。注释中的页码为对应 Markdown 标注的 PDF 页码。
% 880＝27 李林 880 数二高数篇做题本（A4 紧凑版）。
% 500＝李林基础 500 题数二题本（A4 版）。计算＝计算与方法篇，概念＝概念与理解篇。
% 讲义＝27 考研李林辅导讲义（高数）A4 做题本。
% 六套卷＝27 李林考研数学冲刺预测 6 套卷（数二）。
% 26版＝26 版 880 数二高数做题本（A4 版）。
% 部分题目调整了表述或选项顺序，组合题在注释中标明所用小问。
\documentclass[UTF8,a4paper,zihao=-4,fontset=fandol]{ctexart}
\usepackage[margin=20mm]{geometry}
\usepackage{amsmath,amssymb,enumitem,tabularx}
\usepackage[hidelinks]{hyperref}
\setlength{\parindent}{0pt}
\setlist[enumerate]{label=\arabic*.,leftmargin=2em,itemsep=.4em,topsep=.4em}
\setlength{\tabcolsep}{0pt}
\raggedbottom
\newcommand{\blank}{\underline{\hspace{3em}}}
\newcommand{\choices}[4]{%
  \par\noindent\begin{tabularx}{\linewidth}{@{}X@{\hspace{1em}}X@{}}
    (A)\,#1 & (B)\,#2 \\
    (C)\,#3 & (D)\,#4
  \end{tabularx}\par}
\newcommand{\paper}[1]{%
  \pdfbookmark[0]{第#1卷}{paper.\thepage}%
  \begin{center}
    {\LARGE\bfseries 27302 高等数学阶段模拟卷（第#1卷）\par}
    \small 《工科数学分析基础》第三版上册第一章至第二章第四节\par
    满分 150 分\qquad 建议用时 180 分钟
  \end{center}
  \section*{一、选择题}
  第 1--10 题，每题 5 分，共 50 分。每题只有一个选项符合题意。}
\newcommand{\fillsection}{%
  \section*{二、填空题}
  第 11--16 题，每题 5 分，共 30 分。将答案填在题中横线上。}
\newcommand{\solvesection}{%
  \newpage
  \section*{三、解答题}
  第 17 题 10 分，第 18--22 题每题 12 分，共 70 分。写出演算步骤或证明过程。}

\begin{document}
\paper{二}
\begin{enumerate}
% 880/2/第一章基础选择1
\item 函数 $f(x)=|x\sin x|\mathrm{e}^{\cos x}$，$x\in\mathbb R$，是（\quad）。
\choices{单调函数}{周期函数}{偶函数}{有界函数}
% 880/2/第一章基础选择4
\item 若 $x\to+\infty$ 时 $|f(x)|$ 与 $|g(x)|$ 均趋于 $+\infty$，则此时（\quad）。
\choices{$f-g$ 为无穷小}{$f+g$ 为无穷大}{$g/f\to1$}{$(f+g)/(fg)$ 为无穷小}
% 500/3/第一章计算6
\item 设 $c>0$，$f(x)=\begin{cases}x^2-7,&|x|\le c,\\-6/|x|,&|x|>c.\end{cases}$
使 $f$ 在 $\mathbb R$ 上连续的 $c$ 有（\quad）。
\choices{1 个}{2 个}{3 个}{4 个}
% 880/7/第一章综合选择9
\item 下列关于数列的结论中，错误的是（\quad）。
\choices{若 $a_n\to a>1$，则存在 $M>1$，使充分大的 $n$ 满足 $a_n>M$}
{若 $a_n\to a$、$b_n\to b$ 且 $a<b$，则充分大的 $n$ 满足 $a_n<b_n$}
{若 $a_n\to a$ 且对所有 $n$ 有 $M\le a_n\le N$，则 $M\le a\le N$}
{若 $a_n\to a\ne0$，则充分大的 $n$ 满足 $a_n>a-1/n$}
% 500/15/第一章概念8
\item 对有界数列 $\{x_n\},\{y_n\}$，二者至少有一个趋于 $0$，是 $x_ny_n\to0$ 的（\quad）。
\choices{充分非必要条件}{必要非充分条件}{充要条件}{既非充分也非必要条件}
% 880/14/第二章基础选择1
\item 设 $\varphi$ 有界，$f(x)=\begin{cases}(1-\cos x)/x,&x>0,\\x^2\varphi(x),&x\le0.\end{cases}$
则 $f$ 在 $x=0$ 处（\quad）。
\choices{可导}{连续但不可导}{极限存在但不连续}{极限不存在}
% 880/14/第二章基础选择2
\item 设 $f'(x)$ 存在，$a,b\in\mathbb R$。则 $\displaystyle
\lim_{h\to0}\frac{f(x+ah)-f(x-bh)}{h}=$（\quad）。
\choices{$(a+b)f'(x)$}{$(a-b)f'(x)$}{$af'(x)$}{$bf'(x)$}
\end{enumerate}
\newpage
\begin{enumerate}[start=8]
% 500/36/第二章概念5
\item 设 $f$ 在 $a$ 处可导。$|f(x)|$ 在 $a$ 处不可导的充分条件是（\quad）。
\choices{$f(a)=f'(a)=0$}{$f(a)=0,\ f'(a)\ne0$}{$f(a)>0,\ f'(a)>0$}{$f(a)<0,\ f'(a)<0$}
% 880/15/第二章基础选择10
\item 设 $f$ 在 $0$ 的邻域内连续，$f(0)=0$ 且
$\displaystyle\lim_{x\to0}\frac{f(x)}{1-\cos x}=2$，则 $f$ 在 $0$ 处（\quad）。
\choices{不可导}{可导且导数非零}{有极小值}{有极大值}
% 500/37/第二章概念7
\item 下列关于函数可导性的命题中，正确的是（\quad）。
\choices{若 $f,g$ 在 $x_0$ 处均不可导，则 $f+g$ 在 $x_0$ 处必不可导}
{若 $f,g$ 在 $x_0$ 处均不可导，则 $fg$ 在 $x_0$ 处必不可导}
{若 $g$ 在 $x_0$ 处可导，$f$ 在 $g(x_0)$ 处不可导，则 $f\circ g$ 在 $x_0$ 处可能可导}
{若 $g$ 在 $x_0$ 处不可导，$f$ 在 $g(x_0)$ 处可导，则 $f\circ g$ 在 $x_0$ 处必不可导}
\end{enumerate}
\fillsection
\begin{enumerate}[start=11]
% 500/8/第一章计算19
\item $\displaystyle\lim_{n\to\infty}\frac{1\cdot1!+2\cdot2!+\cdots+n\cdot n!}{(n+1)!}=\blank$。
% 880/4/第一章基础填空3
\item 若 $f(x)=\begin{cases}(\sin 2x+\mathrm{e}^{2ax}-1)/x,&x\ne0,\\a,&x=0\end{cases}$
在 $0$ 处连续，则 $a=\blank$。
% 880/4/第一章基础填空4
\item 设 $a>0,\ a\ne1$。使 $\displaystyle\lim_{x\to+\infty}x^p(a^{1/x}-a^{1/(x+1)})$
存在且为有限值的 $p$ 的取值范围为 $\blank$。
% 500/24/第二章计算22
\item 若 $f(x)=\sin^2x\cdot\mathrm{e}^x\cos x\cdot\ln(2+x)$，则 $f''(0)=\blank$。
% 500/30/第二章计算42
\item 设 $f$ 是 $g(x)=\mathrm{e}^x+\ln x$（$x>0$）的反函数，
则 $y=f(x)$ 在横坐标为 $\mathrm{e}$ 的点处的切线为 $\blank$。
% 500/29/第二章计算37
\item 设 $f(x)=\begin{cases}\mathrm{e}^x-1,&x<1,\\x^2+(\mathrm{e}-2)x,&x\ge1.\end{cases}$
在 $[0,3]$ 上满足拉格朗日中值定理的中值点 $\xi=\blank$。
\end{enumerate}
\solvesection
\begin{enumerate}[start=17]
% 880/5/第一章基础解答3(I)(II)(III)
\item \textbf{（10 分）} 求下列极限。
\[
 \text{(1)}\ \lim_{n\to\infty}\sum_{k=1}^n\frac{k}{n^2+n+k},
 \quad
 \text{(2)}\ \lim_{n\to\infty}
 \left(\sqrt{\sum_{k=1}^n k}-\sqrt{\sum_{k=1}^{n-1}k}\right),
 \quad
 \text{(3)}\ \lim_{n\to\infty}\sum_{k=1}^n\frac{1}{4k^2-1}.
\]
\vspace{4.5cm}
% 880/21/第二章基础解答3
\item \textbf{（12 分）} 求各题中的 $\mathrm{d}y$。
\begin{enumerate}[label=(\arabic*),nosep]
\item $y=\varphi(\arctan(1/x))$，其中 $\varphi$ 可导。
\item $y$ 由 $\mathrm{e}^{x+y}=y\sin x$ 确定，在 $\mathrm{e}^{x+y}-\sin x\ne0$ 处求微分。
\item $x=2t,\ y=5t^2+1$，将结果写成关于 $\mathrm{d}x$ 的式子。
\end{enumerate}
\vspace{4.5cm}
% 880/30/第二章综合解答1
\item \textbf{（12 分）} 设 $f(x)=\begin{cases}ax^2+b\sin x+c,&x\le0,\\\ln(1+x),&x>0.\end{cases}$
求使 $f'$ 在 $0$ 处连续、但 $f''(0)$ 不存在的 $a,b,c$。
\end{enumerate}
\newpage
\begin{enumerate}[start=20]
% 讲义/30/第一章连续性例8；讲义/69/第三章中值定理例1
\item \textbf{（12 分）} 证明下列结论。
\begin{enumerate}[label=(\arabic*),nosep]
\item 若 $f$ 在 $[0,3]$ 上连续，$f(0)+f(1)+f(2)=3$ 且 $f(3)=1$，
则存在 $\xi\in[0,2]$ 使 $f(\xi)=f(3)$。
\item 若 $g$ 在 $[0,1]$ 上二阶可导，曲线 $y=g(x)$ 与连接两端点的弦在区间内部相交，
则存在 $\eta\in(0,1)$ 使 $g''(\eta)=0$。
\end{enumerate}
\vspace{5.5cm}
% 880/22/第二章基础解答14
\item \textbf{（12 分）} 设 $0<a<b$，$f$ 在 $[a,b]$ 上连续、在 $(a,b)$ 内可导，且 $f(a)=f(b)=0$。
证明存在 $\xi,\eta\in(a,b)$，分别使
$2f(\xi)+\xi f'(\xi)=0$ 和 $2\eta f(\eta)-f'(\eta)=0$。
\vspace{5.5cm}
% 880/12/第一章综合解答12
\item \textbf{（12 分）} 设 $0<x_1<\pi$，$x_{n+1}=\sin x_n$。
证明 $\{x_n\}$ 收敛并求其极限，再求
$\displaystyle\lim_{n\to\infty}(x_{n+1}/x_n)^{1/x_n^2}$。
\end{enumerate}

\newpage
\paper{三}
\begin{enumerate}
% 880/2/第一章基础选择2
\item 在 $(0,\pi/2)$ 上设 $f(x)=\cos(\sin x)$，$g(x)=\sin(\cos x)$，则（\quad）。
\choices{$f,g$ 均递减}{$f$ 递减，$g$ 递增}{$f,g$ 均递增}{$f$ 递增，$g$ 递减}
% 880/2/第一章基础选择5
\item 当 $x\to0$ 时，$\dfrac{1}{x^2}\sin\dfrac1x$ 是（\quad）。
\choices{无穷大}{无穷小}{有界但非无穷小}{无界但非无穷大}
% 880/3/第一章基础选择8
\item 当 $x\to0$ 时，若 $\mathrm{e}^x-\dfrac{1+ax^2}{1+bx}$ 与 $x^3$ 是同阶无穷小，则（\quad）。
\choices{$a=\tfrac12,\ b=1$}{$a=-\tfrac12,\ b=-1$}{$a=\tfrac12,\ b=-1$}{$a=-\tfrac12,\ b=1$}
% 500/4/第一章计算7
\item 函数 $\displaystyle f(x)=\frac{1/x-1/(x+1)}{1/(x-1)-1/x}$ 的可去间断点个数为（\quad）。
\choices{0}{1}{2}{3}
% 880/7/第一章综合选择10
\item 设 $\{x_n\},\{y_n\}$ 为两个数列。下列结论中，正确的是（\quad）。
\choices{若二者均无界，则 $\{x_n+y_n\}$ 无界}
{若二者均无界，则 $\{x_ny_n\}$ 无界}
{若一者有界、一者无界，则 $\{x_ny_n\}$ 无界}
{若二者均为无穷大量，则 $\{x_ny_n\}$ 为无穷大量}
% 500/18/第二章计算4
\item 设 $x_0\ne0$，$f$ 在 $x_0$ 的邻域内有定义。$f$ 在 $x_0$ 可导是
$\displaystyle\lim_{x\to x_0}\frac{xf(x_0)-x_0f(x)}{x-x_0}$ 存在的（\quad）。
\choices{充分非必要条件}{必要非充分条件}{充要条件}{既非充分也非必要条件}
% 880/15/第二章基础选择13
\item 设 $f$ 有一阶连续导数，$F(x)=f(x)(1+|\sin x|)$。$f(0)=0$ 是 $F$ 在 $0$ 处可导的（\quad）。
\choices{充要条件}{充分非必要条件}{必要非充分条件}{既非充分也非必要条件}
\end{enumerate}
\newpage
\begin{enumerate}[start=8]
% 880/16/第二章基础选择17
\item 设对所有 $x\in(-\delta,\delta)$ 有 $|f(x)|\le x^2$，其中 $\delta>0$。则 $0$ 是（\quad）。
\choices{间断点}{连续但不可导点}{可导点且 $f'(0)=0$}{可导点且 $f'(0)\ne0$}
% 500/35/第二章概念2
\item 设 $f$ 在 $0$ 处连续，且 $\displaystyle\lim_{h\to0}\frac{f(h)^2}{h|h|}$ 存在。则（\quad）。
\choices{左导数存在，右导数未必存在}{右导数存在，左导数未必存在}
{左、右导数均存在}{左、右导数均未必存在}
% 880/24/第二章综合选择1
\item 设 $f$ 在 $1$ 的邻域内可导，$f'$ 严格递减，且 $f(1)=f'(1)=1$。
在该邻域内，除 $x=1$ 外，有（\quad）。
\choices{两侧均有 $f(x)<x$}{两侧均有 $f(x)>x$}
{左侧 $f(x)<x$，右侧 $f(x)>x$}{左侧 $f(x)>x$，右侧 $f(x)<x$}
\end{enumerate}
\fillsection
\begin{enumerate}[start=11]
% 880/4/第一章基础填空5
\item $\displaystyle\lim_{x\to+\infty}\frac{x^3+x^2+1}{\mathrm{e}^x+x^3}(\sin x+\cos x)=\blank$。
% 880/10/第一章综合填空7
\item 设 $0<a_1<a_2$，则 $\displaystyle\lim_{n\to\infty}(a_1^{-n}+a_2^{-n})^{1/n}=\blank$。
% 500/8/第一章计算21
\item $\displaystyle\lim_{x\to\pi/2}(\sin x)^{\tan x}=\blank$。
% 880/19/第二章基础填空19
\item $f(x)=\cos|x|+x|x|$ 在 $0$ 处存在的最高阶导数的阶数为 $\blank$。
% 500/24/第二章计算23
\item 设 $f$ 对任意 $x$ 均满足 $f(1+x)=af(x)$，且 $f'(0)=b$，$ab\ne0$。
则 $f'(2)=\blank$。
% 500/33/第二章计算49
\item 若曲线 $y=x^2\ln x$ 和 $y=a/x$（$a\ne0$）相切，则 $a=\blank$。
\end{enumerate}
\solvesection
\begin{enumerate}[start=17]
% 讲义/3/第一章函数例3
\item \textbf{（10 分）} 判断 $f(x)=\ln(x+\sqrt{1+x^2})$ 的奇偶性，并求它的反函数。
\vspace{4.5cm}
% 880/11/第一章综合解答3
\item \textbf{（12 分）} 下面两个递推数列均取 $x_1=1,\ x_2=2$。分别求它们的极限。
\[
  \text{(1)}\quad x_{n+2}=\frac{3x_{n+1}-x_n}{2},
  \qquad
  \text{(2)}\quad x_{n+2}=\frac{x_n+x_{n+1}}{2}.
\]
\vspace{4.5cm}
% 880/21/第二章基础解答8
\item \textbf{（12 分）} 设 $f(x)=\begin{cases}\mathrm{e}^{-1/x^2},&x\ne0,\\0,&x=0.\end{cases}$
证明 $f^{(n)}(0)$ 对任意正整数 $n$ 都存在，并求其值。
\end{enumerate}
\newpage
\begin{enumerate}[start=20]
% 880/22/第二章基础解答16
\item \textbf{（12 分）} 设 $f$ 在 $[0,1]$ 上连续，在 $(0,1)$ 内可导，且 $f(0)=0,\ f(1)=1$。
证明存在 $x_0\in(0,1)$ 使 $f(x_0)=2(1-x_0)$，并证明存在不同的 $\xi,\eta\in(0,1)$，
使 $f'(\xi)(1+f'(\eta))=2$。
\vspace{5.5cm}
% 880/12/第一章综合解答13
\item \textbf{（12 分）} 设 $x_1>0$，$x_{n+1}=\ln(\mathrm{e}^{x_n}-1)-\ln x_n$。
证明 $\{x_n\}$ 收敛并求其极限。
\vspace{5.5cm}
% 880/32/第二章综合解答10
\item \textbf{（12 分）} 设 $f$ 在 $[a,b]$ 上二阶可导，$f(a)=f(b)=0$，
且 $f'_+(a)f'_-(b)>0$。证明在 $(a,b)$ 内存在不同的 $\xi,\eta$，
使 $f'(\xi)=0,\ f''(\eta)=0$。
\end{enumerate}

\newpage
\paper{四}
\begin{enumerate}
% 880/2/第一章基础选择6
\item 若 $\displaystyle\lim_{x\to+\infty}\left(\frac{x^2}{x+1}-ax-b\right)=0$，则（\quad）。
\choices{$a=1,\ b=1$}{$a=-1,\ b=1$}{$a=1,\ b=-1$}{$a=-1,\ b=-1$}
% 880/2/第一章基础选择7
\item 当 $x\to0$ 时，$(x-\sin x)\tan x$ 比 $\ln(1+x^n)$ 高阶，而 $\ln(1+x^n)$ 比 $x^2$ 高阶。
则正整数 $n=$（\quad）。
\choices{4}{3}{2}{1}
% 500/7/第一章计算17
\item 当 $x\to1$ 时，若 $\displaystyle 1-\frac{m}{1+x+x^2+\cdots+x^{m-1}}\sim x-1$，
则正整数 $m=$（\quad）。
\choices{2}{3}{4}{5}
% 500/12/第一章概念1
\item 对每个 $\varepsilon>0$ 都有无穷多个 $n$ 满足 $|x_n-a|<\varepsilon$，
是 $x_n\to a$ 的（\quad）。
\choices{充分非必要条件}{必要非充分条件}{充要条件}{既非充分也非必要条件}
% 500/13/第一章概念4
\item 设正项数列 $\{a_n\}$ 发散，下列数列中必发散的是（\quad）。
\choices{$\{a_n+\sin n\}$}{$\{a_n/(1+a_n)\}$}{$\{\ln(a_n+1)\}$}{$\{a_n+1/a_n\}$}
% 880/15/第二章基础选择8
\item 设 $f$ 为二阶可导的奇函数，且在 $(0,+\infty)$ 内有 $f'(x)>0,\ f''(x)>0$。
则在 $(-\infty,0)$ 内必有（\quad）。
\choices{$f'<0,\ f''<0$}{$f'<0,\ f''>0$}{$f'>0,\ f''<0$}{$f'>0,\ f''>0$}
% 880/16/第二章基础选择16
\item 设 $y=\ln(1-2x)$，则 $y^{(10)}=$（\quad）。
\choices{$-\dfrac{9!}{(1-2x)^{10}}$}{$\dfrac{9!}{(1-2x)^{10}}$}
{$-\dfrac{9!\,2^{10}}{(1-2x)^{10}}$}{$\dfrac{10!\,2^9}{(1-2x)^{10}}$}
\end{enumerate}
\newpage
\begin{enumerate}[start=8]
% 六套卷/21/第五套选择2
\item 设 $f$ 连续，且 $\displaystyle\lim_{x\to1}\frac{f(2x-1)}{x-1}=2$。
则 $\displaystyle\lim_{x\to0}\frac{f((1+\sin x)^2)-f(1+\sin x)}{x}=$（\quad）。
\choices{0}{2}{1}{$-1$}
% 500/40/第二章概念18
\item 设 $f$ 在 $[a,b]$ 上连续，在 $(a,b)$ 内可导。下列结论正确的是（\quad）。
\choices{若 $f'$ 在 $(a,b)$ 有零点，则 $f$ 在 $[a,b]$ 的两个不同点取值相同}
{若 $f(a)\ne f(b)$，则 $f'$ 在 $(a,b)$ 无零点}
{若 $f(a)<f(b)$，则存在 $\xi\in(a,b)$ 使 $f'(\xi)>0$}
{若 $g=|f|$ 且 $g(a)=g(b)$，则 $g'$ 在 $(a,b)$ 有零点}
% 讲义/74/第三章中值定理应用例2
\item 设 $f$ 在 $[0,1]$ 上二阶可导且 $f''(x)>0$。
下列结论正确的是（\quad）。
\choices{$f'(1)>f'(0)>f(1)-f(0)$}{$f'(1)>f(1)-f(0)>f'(0)$}
{$f(1)-f(0)>f'(1)>f'(0)$}{$f'(1)>f(0)-f(1)>f'(0)$}
\end{enumerate}
\fillsection
\begin{enumerate}[start=11]
% 500/9/第一章计算23
\item 当 $x\to0$ 时，若 $1-\cos(1-\cos x)\sim ax^n$，则 $(a,n)=\blank$。
% 500/9/第一章计算24
\item $\displaystyle\lim_{x\to+\infty}\left(\sqrt[6]{x^6+x^5}-\sqrt[6]{x^6-x^5}\right)=\blank$。
% 880/4/第一章基础填空6
\item $\displaystyle\lim_{x\to0}\frac{\mathrm{e}^{x^2}-\mathrm{e}^{2-2\cos x}}{\mathrm{e}^{x^4}-1}=\blank$。
% 880/19/第二章基础填空14
\item 设 $f(x)=x(2x-1)(3x-2)\cdots(100x-99)$，则 $f'(0)=\blank$。
% 880/19/第二章基础填空17
\item 设 $f$ 可导，且 $\displaystyle\lim_{x\to0}\frac{f(1)-f(1-x)}{2x}=-1$。
则曲线 $y=f(x)$ 在 $(1,f(1))$ 处的切线斜率为 $\blank$。
% 500/24/第二章计算24
\item 设对任意 $x$ 有 $f(x)=kf(x+2)$，且 $f(x)=x(x^2-4)$（$0\le x\le2$）。
若 $f$ 在 $0$ 处可导，则 $k=\blank$。
\end{enumerate}
\solvesection
\begin{enumerate}[start=17]
% 880/4–5/第一章基础解答2(I)(VII)(VIII)，末问在原题中重复标作(VII)
\item \textbf{（10 分）} 求下列极限。
\[
 \text{(1)}\ \lim_{x\to+\infty}\frac{x^2-x\sin x}{x^2+x\sin(1/x)},
 \quad
 \text{(2)}\ \lim_{x\to0}(1-x^2)^{1/(1-\sqrt{1-x^2})},
 \quad
 \text{(3)}\ \lim_{x\to0^+}x^{\sin x}.
\]
\vspace{4.5cm}
% 880/21/第二章基础解答9；880/22/第二章基础解答10
\item \textbf{（12 分）} 求下列变化率。
\begin{enumerate}[label=(\arabic*),nosep]
\item 气体以 $100\,\mathrm{cm}^3/\mathrm{s}$ 的速率注入球形气球。设压力不变，
求半径为 $10\,\mathrm{cm}$ 时半径的增长速率。
\item 动点 $P$ 在曲线 $9y=4x^2$ 上运动，坐标轴单位为厘米，横坐标增长速率为
$30\,\mathrm{cm}/\mathrm{s}$。求 $P$ 经过 $(3,4)$ 时其到原点的距离变化率。
\end{enumerate}
\vspace{4.5cm}
% 880/30/第二章综合解答2
\item \textbf{（12 分）} 设 $f$ 是周期为 $5$ 的连续函数，在 $x=1$ 处可导，且当 $x$ 充分接近 $0$ 时有
$f(1+\sin x)-3f(1-\sin x)=8x+o(x)$。
求曲线 $y=f(x)$ 在 $(6,f(6))$ 处的切线方程。
\end{enumerate}
\newpage
\begin{enumerate}[start=20]
% 26版/34/第一章拓展题1
\item \textbf{（12 分）} 设 $f$ 在 $[a,b]$ 上可导，$|f'(x)|<1$ 且 $a<f(x)<b$，
令 $F(x)=(x+f(x))/2$。证明 $F$ 有唯一不动点 $x_*\in(a,b)$，
并证明从任意 $x_0\in[a,b]$ 出发的迭代 $x_{n+1}=F(x_n)$ 都收敛于 $x_*$。
\vspace{5.5cm}
% 讲义/71/第三章中值定理例5
\item \textbf{（12 分）} 设奇函数 $f$ 在 $[-1,1]$ 上二阶可导，且 $f(1)=1$。
证明存在 $\xi\in(0,1)$ 使 $f'(\xi)=1$，又存在 $\eta\in(-1,1)$ 使 $f''(\eta)+f'(\eta)=1$。
\vspace{5.5cm}
% 880/22/第二章基础解答17
\item \textbf{（12 分）} 设 $f$ 在 $[0,1]$ 上二阶可导，$|f''(x)|\le1$，且在 $(0,1)$ 内取得最小值。
证明 $|f'(0)|+|f'(1)|\le1$。
\end{enumerate}

\newpage
\paper{五}
\begin{enumerate}
% 500/2/第一章计算3
\item 函数 $f(x)=x\sin(1/x^2)$（$x\ne0$）（\quad）。
\choices{在 $|x|\to\infty$ 时为无穷大}{在 $x\to0$ 时为无穷大}
{在定义域上有界}{在定义域上无界}
% 500/5/第一章计算12
\item 设 $\displaystyle f(x)=\lim_{n\to\infty}\frac{2x^n-3x^{-n}}{x^n+x^{-n}}\sin\frac1x$（$x\ne0$）。
则 $f$（\quad）。
\choices{仅有两个第一类间断点}{仅有三个第一类间断点}
{有两个第一类间断点和一个第二类间断点}{有一个第一类间断点和两个第二类间断点}
% 880/8/第一章综合选择12
\item 设 $x_n,y_n>0$ 且 $\mathrm{e}^{x_n}=x_n+\mathrm{e}^{y_n}$，$x_n\to0$。
则当 $n\to\infty$ 时（\quad）。
\choices{$y_n$ 比 $x_n$ 高阶}{$x_n$ 比 $y_n$ 高阶}{两者等价}{两者同阶但不等价}
% 500/14/第一章概念7
\item 设 $x\to x_0$。$f,g$ 的极限都不存在，是 $f+g,f-g$ 的极限恰有一个存在的（\quad）。
\choices{充分非必要条件}{必要非充分条件}{充要条件}{既非充分也非必要条件}
% 880/6/第一章综合选择6
\item 设 $f$ 在 $0$ 处可导，$f'(0)\ne0,\ f(0)=0$，
且 $F(x)=\begin{cases}f(x)/x,&x\ne0,\\f(0),&x=0.\end{cases}$
则 $0$ 对于 $F$ 是（\quad）。
\choices{连续点}{第一类间断点}{第二类间断点}{以上均不正确}
% 880/16/第二章基础选择15
\item 设 $f$ 有任意阶导数且 $f'=f^2$，则对整数 $n>3$ 有 $f^{(n)}(x)=$（\quad）。
\choices{$nf^{n+1}(x)$}{$n!f^{n+1}(x)$}{$f^{2n}(x)$}{$n!f^{2n}(x)$}
% 500/19/第二章计算8
\item 设 $f(x)=x$（$x\le1$），且当 $\frac{n+1}{n}<x\le\frac{n}{n-1}$（$n=2,3,\dots$）时
有 $f(x)=\frac{n+1}{n}$。则 $x=1$ 是 $f$ 的（\quad）。
\choices{第一类间断点}{第二类间断点}{连续但不可导点}{可导点}
\end{enumerate}
\newpage
\begin{enumerate}[start=8]
% 880/25/第二章综合选择7
\item 设 $f$ 在 $(-1,1)$ 内可导，且 $\displaystyle\lim_{x\to0}\frac{f(x)}{x^2}=1$。
下列结论中必成立的是（\quad）。
\choices{$\lim_{x\to0}f'(x)/x$ 存在}{$\lim_{x\to0}f'(x)/x$ 不存在}
{$f'(0)=0,\ f''(0)=2$}{$f(0)$ 是极小值}
% 500/41/第二章概念19
\item 设 $f$ 在 $[a,b]$ 上可导，且 $|f'(x)|<1$。下列命题中正确的是（\quad）。
\par (1) 对不同的 $x_1,x_2\in[a,b]$，$|f(x_1)-f(x_2)|<|x_1-x_2|$。
\par (2) 必有 $c\in(a,b)$ 使 $f'(c)=0$。
\par (3) 若 $f(a)=a$，则对 $x\in(a,b]$ 有 $f(x)<x$。
\par (4) 若 $f((a+b)/2)=0$，则对 $x\in[a,b]$ 有 $|f(x)|<(b-a)/2$。
\choices{(1)(3)}{(1)(4)}{(1)(3)(4)}{(2)(3)(4)}
% 500/41/第二章概念20
\item 设 $f$ 在 $[-1,1]$ 上二阶可导且 $f''(x)>0$。
记 $A=f'(1)$、$B=f(1)-f(0)$、$C=f(0)-f(-1)$，则（\quad）。
\choices{$A>C>B$}{$B>A>C$}{$A>B>C$}{$C>B>A$}
\end{enumerate}
\fillsection
\begin{enumerate}[start=11]
% 500/10/第一章计算25
\item 求 $\displaystyle\lim_{x\to+\infty}x\big[(x+2)\ln(x+2)-2(x+1)\ln(x+1)+x\ln x\big]$，
结果为 $\blank$。
% 500/11/第一章计算30
\item 若 $\displaystyle\lim_{x\to0}\frac{x^2 f(x)+\cos x-1}{x^4}=0$，
则 $\displaystyle\lim_{x\to0}\frac{2f(x)-1}{2x^2}=\blank$。
% 880/18/第二章基础填空8
\item 当 $x\to0$ 时，若 $x-\sin x\cos x\sim ax^b$，则 $(a,b)=\blank$。
% 880/18/第二章基础填空9
\item 当 $x\to0$ 时，$\mathrm{e}^x+\ln(1-x)-1$ 与 $x^n$ 同阶，则正整数 $n=\blank$。
% 500/34/第二章计算54
\item 球形细胞的体积以每天 $400\,\mu\mathrm{m}^3$ 的速率增长。
半径为 $10\,\mu\mathrm{m}$ 时，半径的增长速率为 $\blank$ 微米/天。
% 880/30/第二章综合填空11
\item 设 $y=f(x)$ 由 $x=t^2+1,\ y=4t-t^2$（$t\ge0$）确定。
则 $\displaystyle\lim_{n\to\infty}n\big[f((2n+1)/n)-3\big]=\blank$。
\end{enumerate}
\solvesection
\begin{enumerate}[start=17]
% 880/4/第一章基础解答1；讲义/20/第一章数列极限例2
\item \textbf{（10 分）} 解答下列两问。
\begin{enumerate}[label=(\arabic*),nosep]
\item 若对任意 $x\ne0$ 有 $af(x)+bf(1/x)=c/x$，$|a|\ne|b|$，
求 $f(x)$ 并证明 $f$ 为奇函数。
\item 设 $a>0$，求 $\displaystyle\lim_{n\to\infty}\sqrt[n]{1+a^n}$。
\end{enumerate}
\vspace{4.5cm}
% 880/13/第一章拓展题4
\item \textbf{（12 分）} 证明对每个正整数 $n$ 有
$\dfrac1n-\ln(1+\dfrac1n)<\dfrac1{2n^2}$，
并求 $\displaystyle\lim_{n\to\infty}\prod_{k=1}^n(1+k/n^2)$。
\vspace{4.5cm}
% 880/32/第二章综合解答16
\item \textbf{（12 分）} 设 $f(x)=\sum_{k=1}^{n}a_k\sin(kx)$，其中 $a_k$ 为实数。
求 $f'(0)$。若对任意 $x$ 有 $|f(x)|\le|\sin x|$，证明 $|\sum_{k=1}^{n}ka_k|\le1$。
\end{enumerate}
\newpage
\begin{enumerate}[start=20]
% 880/22/第二章基础解答12
\item \textbf{（12 分）} 设 $f$ 在 $0$ 处连续，且
$\displaystyle\lim_{x\to0}\frac{xf(x)-(1+x)^{2x}+1}{x^2}=1$。
证明 $f$ 在 $0$ 处可导，并求 $f'(0)$。
\vspace{5.5cm}
% 880/33/第二章综合解答19
\item \textbf{（12 分）} 设 $f$ 有二阶连续导数，$f(0)=f'(0)=0$，$f''(0)>0$。
记曲线 $y=f(x)$ 在 $(x,f(x))$ 处的切线与 $x$ 轴交点的横坐标为 $u(x)$。
求 $\displaystyle\lim_{x\to0}\frac{x}{u(x)}$。
\vspace{5.5cm}
% 880/33/第二章拓展题2
\item \textbf{（12 分）} 设 $f$ 在 $[a,b]$ 上二阶可导，$f(a)=f(b)=0$，
且 $|f''(x)|\le M$。用中值定理证明对任意 $x\in[a,b]$ 均有
\[
  |f(x)|\le\frac{M(b-a)^2}{8},\qquad |f'(x)|\le\frac{M(b-a)}2 .
\]
\end{enumerate}

\newpage
\pdfbookmark[0]{补充练习}{extra}
\begin{center}
  {\LARGE\bfseries 27302 阶段补充练习\par}
  \small 《工科数学分析基础》第三版上册第一章至第二章第四节
\end{center}
\section*{一、选择题}
\begin{enumerate}
% 500/2/第一章计算1
\item 非零常数 $a,b$ 满足 $\displaystyle\lim_{x\to1}\left(\frac{x}{x-1}-\frac a{\ln x}\right)=b$。
则（\quad）。
\choices{$a=1,\ b=1/2$}{$a=-1,\ b=1/2$}{$a=1,\ b=-1/2$}{$a=-1,\ b=-1/2$}
% 500/2/第一章计算2
\item 下列等式中成立的是（\quad）。
\choices{$\displaystyle\lim_{x\to0^+}(1+1/x)^x=1$}{$\displaystyle\lim_{x\to0^+}(1+1/x)^x=\mathrm e$}
{$\displaystyle\lim_{x\to+\infty}(1-1/x)^x=-\mathrm e$}{$\displaystyle\lim_{x\to+\infty}(\sin x)/x=1$}
% 880/8/第一章综合选择13
\item 对正整数 $n$，设 $u_n$ 为 $x(1+\ln x)=n$ 在 $x\ge1$ 的实根。下列结论中正确的是（\quad）。
\choices{$\{u_n\}$ 收敛}{$\{u_n/n\}$ 发散}{$\{u_n\ln(n)/n\}$ 收敛}{$\{u_n\ln(n)/n\}$ 发散}
% 880/8/第一章综合选择14
\item 设 $\{x_n\}$ 为实数列。判断下列命题。
\par (1) 若 $\{\arctan x_n\}$ 收敛，则 $\{x_n\}$ 收敛。
\par (2) 若 $\{\arctan x_n\}$ 单调，则 $\{x_n\}$ 收敛。
\par (3) 若 $x_n\in[-1,1]$ 且 $\{x_n\}$ 收敛，则 $\{\arcsin x_n\}$ 收敛。
\par (4) 若 $x_n\in[-1,1]$ 且 $\{x_n\}$ 单调，则 $\{\arcsin x_n\}$ 收敛。
\choices{(1)(2) 正确}{(3)(4) 正确}{(1)(3) 正确}{(2)(4) 正确}
% 500/14/第一章概念5
\item 设 $a_n>0$，$b_n=\sum_{k=1}^n a_k$。下列命题中正确的有（\quad）个。
\par (1) 若 $\{a_n\}$ 有界，则 $\{b_n\}$ 有界。
\par (2) 若 $\{b_n\}$ 有界，则 $\{a_n\}$ 有界。
\par (3) 若 $\{a_n\}$ 收敛，则 $\{b_n\}$ 收敛。
\par (4) 若 $\{b_n\}$ 收敛，则 $\{a_n\}$ 收敛。
\choices{0}{1}{2}{3}
% 500/15/第一章概念9
\item 当 $x\to0$ 时，$f(x),g(x),x$ 两两同阶但不等价。下列各量中，必为 $x$ 的高阶无穷小的有（\quad）个。
\[
 \text{(1)}\ f+g,\qquad\text{(2)}\ fg,\qquad
 \text{(3)}\ f^2,\qquad\text{(4)}\ f(x)\sin x\cos(1/x).
\]
\choices{1}{2}{3}{4}
% 500/16/第一章概念11
\item 设 $x_n\to0$。下列说法错误的是（\quad）。
\choices{若 $x_1=-1$，则 $\{x_n\}$ 有最小值}
{若 $x_1x_2=-1$，则 $\{x_n\}$ 有最大值和最小值}
{对任意 $\varepsilon>0$，充分大的 $n$ 都满足 $0<|x_n|<\varepsilon$}
{$\{|x_n|\}$ 必有最大值}
% 500/7/第一章计算18
\item 设 $\displaystyle f(x)=\lim_{n\to\infty}\frac{1+x^{2n+1}}{1+x^{2n}}\cos(\pi x)$。
则 $f$（\quad）。
\choices{在 $1$ 与 $-1$ 处均连续}{在 $1$ 处连续、在 $-1$ 处不连续}
{在 $1$ 与 $-1$ 处均不连续}{在 $1$ 处不连续、在 $-1$ 处连续}
% 六套卷/6/第二套选择1
\item 设 $\displaystyle f(x)=\lim_{n\to\infty}\frac{2x^{2n}-1}{1+x^{2n}}\sin(1/x)$（$x\ne0$）。
函数 $f$ 的第一类间断点有（\quad）个。
\choices{0}{1}{2}{3}
% 六套卷/11/第三套选择2
\item 设 $f$ 在 $1$ 的邻域内连续，且
$\displaystyle\lim_{x\to0}\frac{\ln[f(1+x)+\mathrm{e}^{x^2}]}{x^2}=2$。则 $x=1$ 是（\quad）。
\choices{驻点且为极大值点}{驻点且为极小值点}{不可导点}{可导点但不是驻点}
% 880/14/第二章基础选择5
\item 函数 $f(x)=(x^2+3x+2)|x^3-x|$ 在 $\mathbb R$ 上的不可导点有（\quad）个。
\choices{1}{2}{3}{4}
% 880/16/第二章基础选择14
\item 设 $f$ 在 $a$ 处连续。下列条件中，能推出 $f$ 在 $a$ 处可导的是（\quad）。
\choices{$\displaystyle\lim_{x\to+\infty}|x|\big[f(a+1/|x|)-f(a)\big]$ 存在}
{$\displaystyle\lim_{x\to0}\frac{f(a+x^3)-f(a)}{x^2}$ 存在}
{$\displaystyle\lim_{h\to0}\frac{f(a+h)-f(a-h)}{2h}$ 存在}
{$\displaystyle\lim_{x\to0}\frac{f(a+x^3)-f(a)}{\tan^3 x}$ 存在}
% 500/38/第二章概念10
\item 设 $f$ 在 $(0,1)$ 内可导。$f'$ 在 $(0,1)$ 内有界是 $f$ 在 $(0,1)$ 内有界的（\quad）。
\choices{充分非必要条件}{必要非充分条件}{充要条件}{既非充分也非必要条件}
% 880/28/第二章综合选择28
\item 若 $\displaystyle\lim_{x\to x_0^-}f'(x)=\lim_{x\to x_0^+}f'(x)=1$，则（\quad）。
\choices{$f$ 在 $x_0$ 处必可导且 $f'(x_0)=1$}{$f$ 在 $x_0$ 处必连续但不可导}
{$f$ 在 $x_0$ 处必存在极限但不连续}{$f$ 在 $x_0$ 的某去心左、右邻域内分别单调递增}
% 500/23/第二章计算19
\item 若 $3a^2-5b<0$，则方程 $x^5+2ax^3+3bx+4c=0$ 有（\quad）。
\choices{0 个实根}{1 个实根}{3 个不同实根}{5 个不同实根}
% 880/3/第一章基础选择9
\item 设 $f(x)=(\ln x)^2$、$g(x)=x^2$、$h(x)=\mathrm{e}^x$。当 $x$ 充分大时，（\quad）。
\choices{$f(x)<g(x)<h(x)$}{$g(x)<h(x)<f(x)$}
{$h(x)<g(x)<f(x)$}{$g(x)<f(x)<h(x)$}
% 880/6/第一章综合选择1
\item 使 $\displaystyle\lim_{x\to+\infty}
\frac{\mathrm{e}^{\sin(1/x)}-1}{(1+1/x)^k-(1+1/x)}$ 存在且非零的充要条件是（\quad）。
\choices{$k\ne1$}{$k>1$}{$k>0$}{与 $k$ 无关}
% 880/6/第一章综合选择2
\item 设 $a>0$，$\displaystyle\lim_{x\to+\infty}x^a
\ln\frac{\arctan(x+1)}{\arctan x}=b$。下列哪组 $(a,b)$ 符合等式？
\choices{$a=2,\ b=2/\pi$}{$a=2,\ b=\pi/2$}{$a=1,\ b=2/\pi$}{$a=1,\ b=\pi/2$}
% 880/8/第一章综合选择15
\item 若正项数列 $\{x_n\}$ 对所有 $n$ 都满足 $x_{n+1}<2-1/x_n$，则 $\{x_n\}$（\quad）。
\choices{收敛于 $0$}{收敛于 $1$}{收敛于 $\sqrt 2$}{发散}
% 880/14/第二章基础选择4
\item 设 $f(x)=\begin{cases}a\sqrt{x},&0\le x\le b,\\\ln x,&x>b\end{cases}$（$b>0$），
且 $f$ 在 $(0,+\infty)$ 内可导。则（\quad）。
\choices{$a=2/\mathrm e,\ b=\mathrm e^2$}{$a=2/\mathrm e,\ b=\mathrm e$}
{$a=1/\mathrm e,\ b=\mathrm e^2$}{$a=1/\mathrm e,\ b=\mathrm e$}
% 500/20/第二章计算12
\item 设 $f(x)=\begin{cases}x^{2x},&x>0,\\x\mathrm{e}^x+1,&x\le0.\end{cases}$
则 $x=0$ 是 $f$ 的（\quad）。
\choices{可导点、极值点}{不可导点、极值点}{可导点、非极值点}{不可导点、非极值点}
\end{enumerate}

\section*{二、填空题}
\begin{enumerate}[start=22]
% 880/4/第一章基础填空1
\item 设 $f(x)=\begin{cases}1,&|x|\le1,\\0,&|x|>1.\end{cases}$ 则 $f(f(f(x)))=\blank$。
% 880/4/第一章基础填空2
\item 当 $x\to0$ 时，若 $(1+ax^2)^{1/3}-1\sim\cos x-1$，则 $a=\blank$。
% 880/9/第一章综合填空1
\item 当 $x\to0$ 时，$3x-4\sin x+\sin x\cos x$ 是关于 $x$ 的 $\blank$ 阶无穷小。
% 880/9/第一章综合填空6
\item 若 $k\ne1/2$，则 $\displaystyle\lim_{n\to\infty}
  n\ln\left[\frac{n-2nk+1}{n(1-2k)}\right]=\blank$。
% 880/10/第一章综合填空8
\item 若 $\displaystyle\lim_{x\to+\infty}(\sqrt[3]{1-x^6}-ax^2-b)=0$，则 $(a,b)=\blank$。
% 880/10/第一章综合填空12
\item 设 $y(x)$ 由 $y^2+xy+x^2-x=0$ 确定，且 $y(1)=-1$，
则 $\displaystyle\lim_{x\to1}\frac{(x-1)^2}{1+y(x)}=\blank$。
% 880/10/第一章综合填空13
\item 若 $f''(1)$ 存在且 $f'(1)\ne0$，则 $\displaystyle\lim_{x\to1}
\left[\frac1{f'(1)(x-1)}-\frac1{f(x)-f(1)}\right]=\blank$。
% 880/17/第二章基础填空1
\item 若 $f(x)=\begin{cases}\arctan(1/x),&x>0,\\ax+b,&x\le0\end{cases}$
在 $0$ 处可导，则 $(a,b)=\blank$。
% 880/18/第二章基础填空7
\item 设 $f'(0)$ 存在且 $f(0)=0$。若 $\displaystyle\lim_{x\to0}
\left[1+\frac{1-\cos f(x)}{\sin x}\right]^{1/x}=\mathrm e$，则 $f'(0)$ 的可能取值为 $\blank$。
% 880/18/第二章基础填空11
\item 设 $f_n(x)=n^2\mathrm{e}^{x/n}-(1+n)x$ 在 $x=x_n$ 处有水平切线。
则 $\displaystyle\lim_{n\to\infty}\mathrm{e}^{x_n}=\blank$。
% 500/25/第二章计算26
\item 若 $\displaystyle y=\ln\frac{\sqrt{1+x^2}-1}{\sqrt{1+x^2}+1}$（$x\ne0$），则 $y'=\blank$。
% 500/26/第二章计算29
\item 设 $y\ge1$ 由 $x=y^{y^2}$ 确定，求用 $\mathrm{d}x$ 表示的 $\mathrm{d}y=\blank$。
% 500/27/第二章计算31
\item 对 $x\ne0$ 设 $\displaystyle f(x)=\lim_{t\to0}(1+t\sin x)^{1/(tx)}$，
则 $\displaystyle\lim_{x\to0}f'(x)=\blank$。
% 500/27/第二章计算32
\item 设 $f$ 可导且 $f'(0)=1$，$y(x)$ 由 $y-x\mathrm{e}^{y-1}=1$ 确定。
若 $z=f(\ln y+\sin x)$，则 $\left.\dfrac{\mathrm{d}z}{\mathrm{d}x}\right|_{x=0}=\blank$。
% 500/28/第二章计算34
\item 设 $y=f(x)$ 由 $\ln(x^2+y)=x^3y+\sin x$ 在 $(0,1)$ 的邻域内确定。
则 $\displaystyle\lim_{n\to\infty}n[f(2/n)-1]=\blank$。
% 500/28/第二章计算35
\item 设 $f(x)=\cos(\cos x)+\cos(\sin x)$，则 $f'''(\pi/2)=\blank$。
% 500/30/第二章计算41
\item 曲线 $x^2+xy+y^3=3$ 在 $x=1$ 的对应点处的法线方程为 $\blank$。
% 500/31/第二章计算45
\item 对数螺线 $r=\mathrm{e}^{\theta}$ 在 $\theta=\pi/2$ 处的切线直角坐标方程为 $\blank$。
% 880/30/第二章综合填空10
\item 设 $f$ 是 $g$ 的反函数，$g(1)=2$，$g'(1)=-\sqrt{3}/3$。则
$\displaystyle\lim_{y\to2}\frac{(y-2)[f(y)-f(2)]}{(\ln y-\ln 2)^2}=\blank$。
% 500/10/第一章计算26
\item $\displaystyle\lim_{x\to0}\left(\frac{1}{x(1+x)^{1/x}}-\frac1{x\mathrm e}\right)=\blank$。
% 500/10/第一章计算27
\item $\displaystyle\lim_{n\to\infty}\sum_{k=1}^{n}\frac{k^2}{\sqrt{n^6+kn}}=\blank$。
% 500/11/第一章计算29
\item 若 $\displaystyle\lim_{x\to+\infty}\left(x^{2/3}\sqrt[3]{6-x}-ax-b\right)=0$，则 $ab=\blank$。
% 880/18/第二章基础填空2
\item 设 $f$ 在 $0$ 处可导，且 $f'(0)=2,\ f(0)=0$。则
$\displaystyle\lim_{x\to0}\frac{f(1-\cos x)}{\ln(1+x^2)}=\blank$。
% 500/27/第二章计算33
\item 设 $g(x)=(\mathrm{e}^x+\mathrm{e}^{-x})/2$，对任意实数 $x,y$ 都有
$f(x+y)=f(x)g(y)+f(y)g(x)$，且 $\displaystyle\lim_{x\to0} f(x)/x=1$。
则 $f'(x)=\blank$。
% 500/28/第二章计算36
\item 设 $f(x)=x^2(1+\sin x)$，则 $f^{(5)}(\pi)=\blank$。
% 500/31/第二章计算43
\item 过点 $(3,-1)$ 作参数曲线 $x=\mathrm{e}^t+t,\ y=\mathrm{e}^{-t}-t$（$t\le1$）的切线，
该切线在 $y$ 轴上的截距为 $\blank$。
% 500/31/第二章计算44
\item 双纽线 $r=\sqrt{\cos(2\theta)}$ 在 $\theta=\pi/12$ 处的切线斜率为 $\blank$。
% 880/29/第二章综合填空1
\item 设 $f(x)=\displaystyle\prod_{k=1}^{100}\big[\tan^k(\pi x/4)-k\big]$，
则 $f'(1)=\blank$。
\end{enumerate}

\section*{三、解答题}
\begin{enumerate}[start=50]
% 880/4/第一章基础解答2(II)
\item 设 $a,b,c>0$，求 $\displaystyle\lim_{x\to+\infty}
\left(\frac{a^{1/x}+b^{1/x}+c^{1/x}}3\right)^x$。
% 880/10/第一章综合解答1
\item 设数列 $\{a_n\}$ 从某一项起各项均非零，且 $a_{n+1}/a_n\to q$，$|q|<1$。
证明 $a_n\to0$。
% 880/10/第一章综合解答2
\item 设 $a_k=2^{1/2^k}$，$u_n=\prod_{k=1}^n a_k$。求 $\lim_{n\to\infty}u_n$。
% 880/12/第一章综合解答14
\item 求下列极限。
\[
\begin{aligned}
  &\text{(1)}\ \lim_{n\to\infty}\prod_{k=1}^n(1+x^{2^k})\quad (|x|<1),\\
  &\text{(2)}\ \lim_{n\to\infty}\prod_{k=1}^n\cos(x/2^k)\quad (x\ne0),\\
  &\text{(3)}\ \lim_{x\to\pi/2}\frac{\prod_{k=2}^{n}(1-\sin^k x)}
    {(1-\sin x)^{n-1}}\quad (n\ge2).
\end{aligned}
\]
% 880/12/第一章综合解答15(II)
\item 设 $a\ne0$，$f$ 为三次多项式，且
$\displaystyle\lim_{x\to2a}\frac{f(x)}{x-2a}=\lim_{x\to4a}\frac{f(x)}{x-4a}=1$。
求 $\displaystyle\lim_{x\to3a}\frac{f(x)}{x-3a}$。
% 880/12/第一章综合解答16
\item 设 $x_1=1/2$，$x_{n+1}=x_n+x_n^2$，求 $\displaystyle\lim_{n\to\infty}\sum_{k=1}^n\frac1{x_k+1}$。
% 880/13/第一章拓展题1(I)
\item 设 $x_0\in(-\pi/2,\pi/2)$，$x_{n+1}=(x_n+\arctan x_n)/2$。证明数列收敛并求极限。
% 880/13/第一章拓展题3(I)
\item 对每个正整数 $n$，证明 $x^n-\cos x=0$ 在 $(0,1)$ 有唯一实根。
% 26版/36/第一章拓展题5
\item 设 $f$ 在 $[0,1]$ 上二阶可导，$f'(x)>0$、$f''(x)>0$，且 $f(0)=0$。
任取 $x_1\in(0,1)$，令 $x_{n+1}=x_n-f(x_n)/f'(x_n)$。证明数列收敛并求极限。
% 880/21/第二章基础解答7
\item 设 $f$ 在 $(0,+\infty)$ 上满足 $f(xy)=f(x)+f(y)$ 且 $f'(1)=1$。
证明 $f$ 在 $(0,+\infty)$ 上可导，并求 $f'(x)$。
% 880/22/第二章基础解答11
\item 设 $f$ 二阶可导且 $f(0)=0$、$f'(0)=1$、$f''(0)=2$。求
$\displaystyle\lim_{x\to0}\frac{f(x)-x}{x^2}$。
% 880/22/第二章基础解答13
\item 设 $f(x)=\begin{cases}1+x^2,&0\le x\le1,\\1-x^2,&-1\le x<0.\end{cases}$
验证拉格朗日中值定理的条件，并求满足定理等式的所有 $\xi$。
% 880/22/第二章基础解答15
\item 设 $f$ 在 $[0,1]$ 上可导，$f(0)=0$，$f(1)=1$，且 $f(x)$ 不恒等于 $x$。
证明存在 $\xi\in(0,1)$，使 $f'(\xi)>1$。
% 880/23/第二章基础解答19
\item 设 $0<a<b$，证明至少存在 $\xi\in(a,b)$，使
$\displaystyle\frac{a\mathrm e^b-b\mathrm e^a}{a-b}=\mathrm e^\xi(1-\xi)$。
% 880/23/第二章基础解答20
\item 证明下列不等式。
\begin{enumerate}[label=(\arabic*),nosep]
\item 当 $0<x<\pi$ 时，$\sin(x/2)>x/\pi$。
\item 当 $\mathrm e<a<b$ 时，$a^b>b^a$。
\item 当 $x>0$ 时，$(x^2-1)\ln x\ge(x-1)^2$。
\item 若 $f$ 在 $\mathbb R$ 上二阶可导，$f''(x)>0$ 且 $\lim_{x\to0}f(x)/x=1$，则 $f(x)\ge x$。
\end{enumerate}
% 880/23/第二章基础解答21
\item 证明当 $x>0$ 时，$\big(1+\frac1{2x}\big)\big(1+\frac1x\big)^x>\mathrm e$。
% 880/31/第二章综合解答6
\item 设 $f$ 在 $[0,1]$ 上二阶可导，且
$\displaystyle\lim_{x\to0^+}\frac{f(x)}x=\lim_{x\to1^-}\frac{f(x)}{x-1}=1$。
证明在 $(0,1)$ 内存在 $\xi,\eta$，使 $f'(\xi)=0$、$f''(\eta)=f(\eta)$。
% 880/32/第二章综合解答12(II)
\item 证明对任意 $x\ge1$，$\arctan x-\frac12\arccos\frac{2x}{1+x^2}=\frac\pi4$。
% 880/32/第二章综合解答15
\item 设 $f(x)=\arctan x$，求 $f^{(n)}(0)$（$n$ 为正整数）。
% 880/32/第二章综合解答17
\item 设 $f$ 可导且 $f(x)>0$，证明曲线 $y=f(x)$ 与 $y=f(x)\sin x$ 在交点处相切。
% 880/33/第二章综合解答20
\item 设 $f$ 在 $(a,b)$ 内可导且 $f'$ 严格递增。
用中值定理证明对不同的 $x_1,x_2\in(a,b)$ 和 $0<\lambda<1$ 有
$f(\lambda x_1+(1-\lambda)x_2)<\lambda f(x_1)+(1-\lambda)f(x_2)$。
% 六套卷/6/第二套选择2，改为证明极限
\item 设 $f$ 在 $[0,+\infty)$ 上有一阶连续导数且 $f(x)>0$，
$\displaystyle\lim_{x\to+\infty}[f(x)+f'(x)]=0$。证明 $\displaystyle\lim_{x\to+\infty}f(x)=0$。
% 880/5/第一章基础解答2(III)(IV)
\item 求下列极限。
\[
 \text{(1)}\quad \lim_{x\to0}\frac{\ln(\sin^2x+\mathrm e^x)-x}{\ln(\mathrm e^{2x}-x^2)-2x},
 \qquad
 \text{(2)}\quad \lim_{x\to0}\frac{(1+x)^3-\mathrm e^{3x}}{x^2}.
\]
% 880/5/第一章基础解答5
\item 对 $x\ne0$，设 $\displaystyle f(x)=\lim_{n\to\infty}
\frac{x^{n+2}-x^{-n}}{x^n+x^{-n}}$。讨论 $f$ 的连续性，并指出间断点的类型。
% 880/20/第二章基础解答1(III)(IV)
\item 分别求 $y=2^{|\sin x|}$ 在 $x\notin\pi\mathbb{Z}$ 处的导数和
$y=\ln|\tan x+\sec x|$ 在其定义域内的导数。
% 880/21/第二章基础解答6
\item 求心形线 $r=1-\cos\theta$ 在 $\theta=\pi/2$ 处的切线直角坐标方程。
% 讲义/54/第二章反函数求导例8
\item 已知 $f'(x)=a\mathrm{e}^x$（$a>0$），求 $f$ 的反函数在 $y=f(x)$ 处的二阶导数。
\end{enumerate}
\end{document}
